%----------------------------------------------------------------------------------------
% Cover Letter Template — Ruilin (Peter) Xie
% Usage: Replace all [PLACEHOLDER] blocks before sending.
% Compile with pdflatex on Overleaf.
%----------------------------------------------------------------------------------------
\documentclass[letterpaper,11pt]{article}

\usepackage[margin=1in]{geometry}
\usepackage{hyperref}
\usepackage{fontenc}
\usepackage{inputenc}
\usepackage{parskip}
\usepackage{microtype}

\hypersetup{colorlinks=true, urlcolor=black, pdfborder={0 0 0}}
\setlength{\parindent}{0pt}
\setlength{\parskip}{10pt}

%----------------------------------------------------------------------------------------
\begin{document}
\pagestyle{empty}

% ── Header ──────────────────────────────────────────────────────────────────
Ruilin (Peter) Xie\\
ruilinx.peter@gmail.com\ \ $\cdot$\ \ +1\,(626)\,757-9218\\
\href{https://misakarinowo.github.io/Portfolio}{misakarinowo.github.io/Portfolio}\ \ $\cdot$\ \
\href{https://github.com/MisakaRinOwO}{github.com/MisakaRinOwO}

\today

% ── Salutation ───────────────────────────────────────────────────────────────
% If you have a specific name, use "Dear [Name],"
% Otherwise use the line below.
Dear [COMPANY NAME] Team,

% ── Paragraph 1: Why this company (most important — must be specific) ────────
% Replace this entire paragraph. 1–3 sentences.
% What draws you to them: a specific game mechanic, design philosophy,
% a system you found interesting as a player, or a public talk/blog post.
% Do NOT write "I have always been passionate about games."
[WHY THIS COMPANY: e.g. "What draws me to [Studio] is [specific system / design decision
in [Game]] — the way [brief observation]. That kind of design — [one-sentence takeaway] —
is exactly the problem space I want to work in."]

% ── Paragraph 2: Who you are + strongest evidence (keep concise) ─────────────
I'm a gameplay programmer and technical designer with a focus on runtime systems and
data-driven architecture. My recent work includes:

\begin{itemize}
  \setlength\itemsep{2pt}
  \item \textbf{CrossBolt} (Unity, C\#): owned the full chart pipeline from authoring grammar
    to runtime spawner and a mainline judgement-path optimization
    (O(N)$\rightarrow$O(P) pending-window). Code is public on GitHub.
  \item \textbf{Void Architect} (UE5, Blueprint): solo modular ship-combat prototype with a
    tick-driven simulation, per-tick replay UI, and Python-assisted balance iteration.
  \item \textbf{Project Summon} (UE5, Blueprint/C++): team action RPG shipped on a fixed
    10-week schedule; owned room progression, combat, inventory, and multiplayer
    persistence systems.
\end{itemize}

% ── Paragraph 3: Role fit + what you bring (tailor to GP vs TD) ──────────────
% GP version:
I'm applying for the [ROLE TITLE] position. My background covers end-to-end runtime
system ownership — from data schema design through spawning, hit detection, and
optimization — alongside experience shipping under real schedule constraints. I'm
comfortable working across Blueprint and C++, and I'm used to integrating with
teammate-owned systems rather than working in isolation.

% TD version (swap in if applying to a TD role):
% I'm applying for the [ROLE TITLE] position. My background sits at the intersection of
% gameplay engineering and systems design: I design data schemas that encode player-facing
% intent, write the runtime code that executes them, and iterate on balance with tooling
% I build myself. I'm most effective when design decisions need to be grounded in
% implementation constraints, and vice versa.

% ── Paragraph 4: Close ───────────────────────────────────────────────────────
My full portfolio is at \href{https://misakarinowo.github.io/Portfolio}{misakarinowo.github.io/Portfolio}.
I'd welcome the chance to talk about how my work fits what your team is building.

Thank you for your time.

Sincerely,\\[6pt]
Ruilin (Peter) Xie

\end{document}
